\section{Boundary Loss Implementation}
\label{sec:boundary_loss}
We explicitly handle boundaries:  
(i) \textbf{Periodic domains (Heat, SWE)}: soft equality between opposing faces, optionally including derivatives:
\[
\mathcal{L}_{\text{bc}}^{\text{per}} = \| u(x=0,y,t) - u(x=L_x,y,t) \|_2^2.
\]  
(ii) \textbf{Open domains (Pollution)}: a sponge rim loss damps edges, and an Orlanski-style radiation condition enforces
\[
u_t + c_{\text{eff}} \,\partial_n u \approx 0,
\]
with $c_{\text{eff}}$ estimated from gradients and clamped for stability. Both losses are formulated in normalized coordinates and evaluated via autograd.\\

\section{Formal Definitions and Assumptions}
\label{sec:definitions}
\begin{definition}[Finite-order local differential operator]
\label{def:local_operator}
Let $u:\mathbb{R}^{d+1}\to\mathbb{R}^m$.  
A differential operator $\mathcal{L}$ is called a \emph{finite-order local operator of order $r$} if
\[
\mathcal{L}u(z) \;=\; F\!\left(z,\, \{D^\alpha u(z) : |\alpha|\le r\}\right),
\qquad z\in\mathbb{R}^{d+1},
\]
for some continuous function $F$, where $\alpha$ is a multi-index.  

\begin{itemize}[topsep=0pt,leftmargin=*]
    \item \textbf{Finite-order:} Only derivatives of $u$ up to order $r<\infty$ appear.
    \item \textbf{Local:} $\mathcal{L}u(z)$ depends only on $u$ and its derivatives evaluated at the same point $z$, not on integrals or values of $u$ away from $z$.
\end{itemize}

\noindent Examples include the heat operator $\partial_t - \Delta$, the wave operator $\partial_{tt}-c^2\Delta$, and the advection operator $\partial_t+v\cdot\nabla$.  
Nonlocal operators such as the fractional Laplacian $(-\Delta)^{\alpha/2}$ are excluded.
\end{definition}

\begin{assumption}
\label{ass:locality}
\noindent\textbf{Locality Assumption:} Assume the setup of Theorem~\ref{prop:universal_pde} with stencil
$\mathcal{S}\subset \{-s_{\max},\ldots,s_{\max}\}^d\times\{0,\ldots,k\}$ of size $N:=|\mathcal{S}|$.
Let the explicit update at $(\mathbf{i},n{+}1)$ be
\[
u^{n+1}(\mathbf{i}) \;=\; \psi\!\left(\{u^{n-\ell}(\mathbf{i}+\mathbf{s}) : (\mathbf{s},\ell)\in \mathcal{S}\}\right),\]
\[
\qquad \psi:\mathbb{R}^{N}\to\mathbb{R}^q.
\]
\end{assumption}


% --- Longitudinal (space-only) sampling model and bound ---
\begin{assumption}
\label{ass:sampling}
\noindent\textbf{Sampling Model:}
Let the stencil factor as
\[
\mathcal{S} = \mathcal{S}_x \times \mathcal{T},
\qquad \mathcal{S}_x \subset \{-s_{\max},\ldots,s_{\max}\}^d,\]
\[|\mathcal{S}_x|=N_x,
\qquad \mathcal{T}=\{0,1,\ldots,k\}.
\]
Thus the stencil size is $N=N_x(k+1)$, but coverage is determined entirely by the $N_x$ spatial
indices. A sensor placed at spatial offset $s\in\mathcal{S}_x$ provides all $k{+}1$ time levels.

We randomly place $m_x$ sensors \emph{uniformly without replacement} among the $N_x$ essential
spatial locations. FieldFormer then takes all available measurements.
\end{assumption}

\begin{assumption}
\label{ass:influence}
\noindent\textbf{Influence Margin (Average-Case):}
Let $V\sim\mathcal{D}$ be the random input (stencil values), and $\psi:\mathbb{R}^{N}\to\mathbb{R}^{q}$ the local update map.
Index the stencil as $\mathcal{S}=\mathcal{S}_x\times\mathcal{T}$ with $\mathcal{S}_x\subset\{-s_{\max},\ldots,s_{\max}\}^d$ and $\mathcal{T}=\{0,\ldots,k\}$, so $N=|\mathcal{S}|=N_x(k{+}1)$ with $N_x=|\mathcal{S}_x|$.

Assume there exist nonnegative weights $(a_{(s,\ell)})_{(s,\ell)\in\mathcal{S}}$ with $$\sum_{(s,\ell)\in\mathcal{S}} a_{(s,\ell)}=1$$ and a constant $c>0$ such that
for every coordinate $(s,\ell)\in\mathcal{S}$ there is a perturbation mechanism $T_{(s,\ell)}$ acting only on $(s,\ell)$ with
\begin{equation}
  \mathbb{E}\bigl[ \,\|\psi(V)-\psi(T_{(s,\ell)} V)\|_2 \,\bigr] \;\ge\; c\, a_{(s,\ell)} .
\label{eq:avg_inf}
\end{equation}
For longitudinal sampling (a spatial miss hides all time levels), aggregate the weights per spatial offset
\[
\bar a_s \;:=\; \sum_{\ell\in\mathcal{T}} a_{(s,\ell)}, \qquad s\in\mathcal{S}_x,
\]
so that $\bar a_s\ge 0$ and $\sum_{s\in\mathcal{S}_x}\bar a_s=1$. For a subset $J\subseteq\mathcal{S}_x$, define its \emph{influence mass} as $A(J):=\sum_{s\in J}\bar a_s$. 
\end{assumption}

\section{Proofs of Theorems}
\label{sec:proofs}

\subsection{Proof of Theorem \ref{prop:universal_pde}}
\begin{theorem}[Expressivity for Locally Structured Spatio-Temporal Dynamics]

Let $u:\mathbb{R}^{d}\times\mathbb{R}\to\mathbb{R}^q$ solve a parabolic or hyperbolic PDE
\[
F\bigl(z,u(z),\nabla u(z),\ldots,D^r u(z)\bigr)=0,\qquad z=(x,t),
\]
with a finite-order, local operator. Assume a consistent \emph{explicit} finite-difference scheme
with compact stencil $\mathcal{S}\subset\{-s_{\max},\ldots,s_{\max}\}^d\times\{0,\ldots,k\}$ and a
CFL-admissible time step $\Delta t$. Then for any $\varepsilon>0$ and compact $K\subset\mathbb{R}^{d+1}$,
there exists a FieldFormer with neighborhood size $m\ge|\mathcal{S}|$, hidden width $d'$, and depth $L$ such that
\[
\sup_{z\in K}\|u(z)-\hat u(z)\|_2<\varepsilon.
\]
\end{theorem}


\begin{proof}[Proof]
Let $u:\mathbb{R}^{d}\times\mathbb{R}\to\mathbb{R}^q$ solve a parabolic or hyperbolic PDE
\[
F\bigl(z,u(z),\nabla u(z),\ldots,D^{r}u(z)\bigr) \;=\; 0,\qquad z=(x,t),
\]
whose differential operator is finite-order and local (Def. 3.1).
Assume a consistent explicit finite-difference scheme with compact stencil
$\mathcal{S}\subset\{-s_{\max},\ldots,s_{\max}\}^d\times\{0,\ldots,k\}$ and CFL-admissible $\Delta t>0$.

\paragraph{1) Discrete locality and finite map.}
On a uniform grid with spatial step $h>0$ and time step $\Delta t>0$, the update rule at index-time $(\mathbf{i},n{+}1)$
is given by
\[
u^{n+1}(\mathbf{i}) \;=\;
\psi\!\left(\{\,u^{n-\ell}(\mathbf{i}+\mathbf{s}) : (\mathbf{s},\ell)\in\mathcal{S}\,\}\right),
\]
for a continuous function $\psi:\mathbb{R}^{N}\to\mathbb{R}^{q}$, where $N:=|\mathcal{S}|$.
Continuity follows because $\psi$ consists of finitely many arithmetic operations
on the stencil entries and coefficients of the PDE discretization.

\paragraph{2) Tokenization of the stencil.}
Let $\Xi=\{(\mathbf{s},\ell):(\mathbf{s},\ell)\in\mathcal{S}\}$ and fix an ordering
$\pi:\{1,\ldots,N\}\to\Xi$ (e.g., lexicographic in $\ell$ then $\mathbf{s}$).
A FieldFormer uses $m$ tokens ($m\ge N$); the first $N$ encode the stencil
values and their relative coordinates. Each token concatenates
(i) relative position $(h\,\mathbf{s}_j,\,\Delta t\,\ell_j)$ and
(ii) the field value $u^{n-\ell_j}(\mathbf{i}+\mathbf{s}_j)$,
producing an input vector $x\in\mathbb{R}^{m d_{\mathrm{in}}}$.
If $m>N$, padded tokens are masked so they have no effect.
Hence there exists a continuous encoder $E:\mathbb{R}^{N}\to\mathbb{R}^{m d_{\mathrm{in}}}$.

\paragraph{3) Compact domain of interest.}
Fix a compact domain $K\subset\mathbb{R}^{d+1}$.
The set of all stencil value tuples attained over $K$,
\[
\mathcal{V}_K
\;:=\;
\Bigl\{
  \bigl(u^{\,n-\ell_j}(\mathbf{i}+\mathbf{s}_j)\bigr)_{j=1}^{N}
  \;:\;
  (x_{\mathbf{i}},\,t_{n+1}) \in K
\Bigr\},
\qquad (x_{\mathbf{i}},t_n) = (\mathbf{i}h,\, n\Delta t).
\]
is compact in $\mathbb{R}^{N}$. The local update
$f:=\psi$ is continuous on $\mathcal{V}_K$, hence $f\circ E^{-1}$ is continuous
on the compact set $E(\mathcal{V}_K)$.


\paragraph{4) Universal approximation by the FieldFormer.}
A transformer/MLP stack defines a continuous function
$g_\theta:\mathbb{R}^{m d_{\mathrm{in}}}\to\mathbb{R}^{q}$.
By universal approximation theorems for MLPs and transformers
\citep{hornik1991approximation,yun2019transformers},
for any $\varepsilon>0$ there exist width $d'$ and depth $L$ such that
\[
\sup_{x\in E(\mathcal{V}_K)}
\|\,g_\theta(x) - f(E^{-1}(x))\,\|_2 < \varepsilon.
\]
Because $E$ is fixed and continuous,
\[
\sup_{v\in\mathcal{V}_K} \|\,g_\theta(E(v)) - \psi(v)\,\|_2 < \varepsilon.
\]
Evaluating $g_\theta(E(v))$ corresponds to the FieldFormer’s prediction
at $(\mathbf{i},n{+}1)$, so each grid update on $K$
is approximated within $\varepsilon$.

\paragraph{5) Uniformity on compact sets.}
Since $K$ is compact and covered by finitely many such neighborhoods,
the uniform bound extends over all of $K$:
\[
\sup_{z\in K}\|u(z)-\hat u(z)\|_2 < \varepsilon.
\]
Hence for some $m\ge|\mathcal{S}|$, width $d'$, and depth $L$,
FieldFormer approximates $u$ uniformly on $K$, completing the proof.
\end{proof}

\subsection{Proof of Theorem \ref{prop:avg_long}}
\begin{theorem}[Average-Case Lower Bound Under Longitudinal Sampling]
Sample $m_x$ spatial locations uniformly without replacement from $\mathcal{S}_x$, let $I\subseteq\mathcal{S}_x$ be the sensed set and $I^c$ the missed set.
Under equation \eqref{eq:avg_inf}, FieldFormer $\hat u$, using all sensed data, satisfies
\[
\inf_{\hat u}\;\mathbb{E}\bigl[\|\hat u-u^{n+1}(\mathbf{i})\|_2\bigr]
\;\;\ge\;\;
c\;\mathbb{E}\bigl[\,A(I^c)\,\bigr]
\;=\;
c\,\Bigl(1-\frac{m_x}{N_x}\Bigr),
\]
where the expectation is over the random sampling of $I$ and the data distribution $\mathcal{D}$.
\end{theorem}

\begin{proof}[Proof]
Let $V\in\mathbb{R}^{N}$ denote the random vector of stencil values
at $(\mathbf{i},n{+}1)$ under distribution $\mathcal{D}$,
with $\mathcal{S}=\mathcal{S}_x\times\mathcal{T}$ and $N=|\mathcal{S}|$.
The target is $\psi(V)\in\mathbb{R}^q$.

\paragraph{1) Conditioning on sensed sites.}
Let $I\subseteq\mathcal{S}_x$ be the set of $m_x$ observed spatial offsets;
for each $s\in I$ all $k{+}1$ time coordinates $(s,\ell)$ are known,
while coordinates with $s\in I^c$ are unobserved.
Any estimator using all available data must be a function
$\hat u=\hat u(V_{I\times\mathcal{T}})$, independent of unobserved values.

\paragraph{2) Indistinguishability and symmetrization.}
Fix $I$. Consider any random transformation $T$ that acts only on
the unobserved coordinates $I^c\times\mathcal{T}$ and leaves
the law of the observed block $V_{I\times\mathcal{T}}$ unchanged.
Then $(V,TV)$ are conditionally indistinguishable to $\hat u$,
meaning $\hat u(V)=\hat u(TV)$ given $I$.
Applying the triangle inequality for the Euclidean norm to the triple
$(a,b,c)=(\psi(V),\,\hat u(V),\,\psi(TV))$ yields
\[
\|\psi(V)-\psi(TV)\|_2
\;\le\;
\|\psi(V)-\hat u(V)\|_2
+ \|\hat u(TV)-\psi(TV)\|_2.
\]
Taking the conditional expectation with respect to $I$ gives
\begin{equation}
\begin{aligned}
\mathbb{E}\!\Big[
  \|\hat u(V)-\psi(V)\|_2
  + \|\hat u(TV)-\psi(TV)\|_2
\Big]
&\;\ge\;
\mathbb{E}\!\Big[
  \|\psi(V)-\psi(TV)\|_2
\Big] \\
&\qquad\text{conditioned on } I.
\end{aligned}
\tag{1}
\end{equation}


Because $\hat u$ depends only on the observed coordinates,
we cannot assume any symmetry between $\psi(V)$ and $\psi(TV)$.
To handle this, we introduce a simple symmetrization.
Let $B\sim\mathrm{Bernoulli}(1/2)$ be an independent coin
and define a randomized input
\[
W \;:=\;
\begin{cases}
V, & B=0,\\[2pt]
TV, & B=1.
\end{cases}
\]
Conditioning on $I$ and averaging (1) over the coin $B$ gives
\begin{equation}
\begin{aligned}
\mathbb{E}\!\left[\|\hat u(W)-\psi(W)\|_2\right]
&=\tfrac12\,\mathbb{E}\!\left[
  \|\hat u(V)-\psi(V)\|_2
  + \|\hat u(TV)-\psi(TV)\|_2
\right] \\
&\ge \tfrac12\,\mathbb{E}\!\left[
  \|\psi(V)-\psi(TV)\|_2
\right],
\qquad \text{given } I.
\end{aligned}
\tag{2}
\end{equation}


Inequality (2) shows that, on average over the randomized choice between
$V$ and $TV$, the estimator’s expected error must be at least half
of the expected discrepancy between the two true outputs.
Consequently, for at least one of $V$ or $TV$ the same bound holds,
so that
\[
\mathbb{E}\!\left[\|\hat u(V)-\psi(V)\|_2 \mid I\right]
\;\ge\;
\tfrac12\,\mathbb{E}\!\left[\|\psi(V)-\psi(TV)\|_2 \mid I\right],
\]
where the factor $\tfrac12$ is absorbed into the constant~$c$ later.


\paragraph{3) Expected influence over unobserved coordinates.}
By Assumption 3.6, for each $(s,\ell)\in\mathcal{S}$ there exists
$T_{(s,\ell)}$ acting on that coordinate only such that
\[
\mathbb{E}\bigl[\|\psi(V)-\psi(T_{(s,\ell)}V)\|_2\bigr] \ge c\,a_{(s,\ell)}.
\]
Summing over unobserved spatial sites and their time levels,
\[
\mathbb{E}\!\left[\|\hat u(V)-\psi(V)\|_2 \,\big|\, I\right]
\ge c \sum_{s\in I^c}\sum_{\ell\in\mathcal{T}} a_{(s,\ell)}
= c\,A(I^c),
\]
where $A(I^c):=\sum_{s\in I^c}\bar a_s$ and $\bar a_s=\sum_{\ell\in\mathcal{T}} a_{(s,\ell)}$.

\paragraph{4) Averaging over random sampling.}
For uniform $m_x$-of-$N_x$ sampling without replacement,
each spatial site is missed with probability $1-m_x/N_x$, so
\[
\mathbb{E}\bigl[A(I^c)\bigr]
= \sum_{s\in\mathcal{S}_x}\bar a_s\,\Pr(s\notin I)
= \Bigl(1-\frac{m_x}{N_x}\Bigr)\!\sum_s \bar a_s
= 1-\frac{m_x}{N_x}.
\]
Taking total expectation over $I$ and $V$ yields
\[
\inf_{\hat u}\,\mathbb{E}\bigl[\|\hat u-u^{n+1}(\mathbf{i})\|_2\bigr]
\ge c\,\mathbb{E}\bigl[A(I^c)\bigr]
= c\Bigl(1-\frac{m_x}{N_x}\Bigr),
\]
as claimed.
\end{proof}

% \section{Additional Results}
% \label{sec:additional_results}

% We report additional results that provide deeper diagnostic insight into reconstruction accuracy and physical consistency across benchmarks. These results are included for completeness and to support the qualitative and quantitative claims made in the main text, but are not used as primary evaluation criteria.

% \begin{table*}[t]
% \centering
% \caption{Physics consistency on the shallow-water equations (SWE) dataset. Metrics are reported with multipliers factored out for clarity. We report relative PDE residual errors (test and full).}
% \label{tab:swe_results_physics}
% \small
% % \resizebox{\linewidth}{!}{
% \begin{tabular}{lcccc}
% \toprule
% Metric & FieldFormer & SIREN & Fourier-MLP & SVGP \\
% \midrule
% Rel. Physics RMSE (test)               & 0.782 & 0.169 & 0.190 & 0.932 \\
% Rel. Physics MAE (test)                & 0.716 & 0.0672 & 0.0753 & 0.918 \\
% Rel. Physics RMSE (full)               & 0.781 & 0.167 & 0.184 & 0.932 \\
% Rel. Physics MAE (full)                & 0.716 & 0.0662 & 0.0732 & 0.917 \\
% \bottomrule
% \end{tabular}
% % }
% \end{table*}

% \begin{table*}[t]
% \centering
% \caption{Performance on the periodic heat equation dataset. We report RMSE and MAE on the sensor test set (with bootstrap mean $\pm$ std), full-field reconstruction errors, and PDE residual and relative residual errors (test and full).}
% \label{tab:heat_results}
% \resizebox{\linewidth}{!}{
% \begin{tabular}{lcccc}
% \toprule
% Metric & FieldFormer & SIREN & Fourier-MLP & SVGP \\
% \midrule
% Test RMSE (bootstrap)     & $1.133\times 10^{-2} \pm 6.3\times 10^{-6}$ & $7.983\times 10^{-2} \pm 6.7\times 10^{-5}$ & $1.259\times 10^{-1} \pm 1.61\times 10^{-4}$ & $5.924\times 10^{-3} \pm 3.4\times 10^{-6}$ \\
% Test MAE (bootstrap)      & $7.715\times 10^{-3} \pm 2.8\times 10^{-6}$ & $5.647\times 10^{-2} \pm 2.5\times 10^{-5}$ & $2.532\times 10^{-2} \pm 3.6\times 10^{-5}$ & $3.627\times 10^{-3} \pm 1.6\times 10^{-6}$ \\
% Full-field RMSE           & $1.133\times 10^{-2}$ & $7.978\times 10^{-2}$ & $1.257\times 10^{-1}$ & $5.923\times 10^{-3}$ \\
% Full-field MAE            & $7.714\times 10^{-3}$ & $5.645\times 10^{-2}$ & $2.530\times 10^{-2}$ & $3.628\times 10^{-3}$ \\
% \bottomrule
% \end{tabular}
% }
% \end{table*}

% \begin{table*}[t]
% \centering
% \caption{Physics consistency on the periodic heat equation dataset. We report PDE residual and relative residual errors (test and full).}
% \label{tab:heat_results_physics}
% \begin{tabular}{lcccc}
% \toprule
% Metric & FieldFormer & SIREN & Fourier-MLP & SVGP \\
% \midrule

% Physics RMSE (test)       & $1.839$ & $2.296$ & $0.631$ & $7.657$ \\
% Physics MAE (test)        & $1.374$ & $1.556$ & $0.137$ & $2.978$ \\
% Physics RMSE (full)       & $1.831$ & $2.298$ & $0.612$ & $7.192$ \\
% Physics MAE (full)        & $1.371$ & $1.551$ & $0.131$ & $2.902$ \\
% Rel. Physics RMSE (test)  & $0.660$ & $0.638$ & $0.246$ & $0.674$ \\
% Rel. Physics MAE (test)   & $0.603$ & $0.538$ & $0.071$ & $0.575$ \\
% Rel. Physics RMSE (full)  & $0.658$ & $0.639$ & $0.240$ & $0.674$ \\
% Rel. Physics MAE (full)   & $0.600$ & $0.539$ & $0.068$ & $0.575$ \\
% \bottomrule
% \end{tabular}
% \end{table*}

% \noindent\textbf{Physics Consistency on the Shallow-Water Equations.}
% Table~\ref{tab:swe_results_physics} reports relative PDE residual errors for the shallow-water equations benchmark. These metrics measure the extent to which reconstructed fields satisfy the governing equations when evaluated via automatic differentiation, normalized by characteristic derivative scales. While FieldFormer does not explicitly minimize the PDE residual, it achieves balanced residual values that are substantially lower than SVGP and comparable in magnitude to neural field baselines. SIREN and Fourier-MLP exhibit lower relative residuals in this setting, reflecting their tendency toward smoother solutions that partially satisfy differential constraints despite higher reconstruction error. These results illustrate that low PDE residuals alone do not guarantee accurate field recovery, particularly in oscillatory, multi-variable systems.

% \noindent\textbf{Reconstruction Performance on the Periodic Heat Equation.}
% Table~\ref{tab:heat_results} reports reconstruction accuracy on the periodic heat equation benchmark, including RMSE and MAE on held-out sensors (with bootstrap statistics) as well as full-field errors. FieldFormer achieves substantially lower reconstruction error than neural field baselines, demonstrating effective exploitation of anisotropic structure under sparse spatial sensing. SVGP attains low numerical error due to its strong smoothness bias, but this behavior is further examined through physics diagnostics in Table~\ref{tab:heat_results_physics}. These results complement the diagnostic discussion in the main text, where the heat equation is used to analyze strengths and architectural limitations under fully specified physics.

% \noindent\textbf{Physics Consistency on the Periodic Heat Equation.}
% Table~\ref{tab:heat_results_physics} reports absolute and relative PDE residual errors for the periodic heat equation. Absolute residuals reflect raw violations of the governing PDE, while relative residuals normalize these values to account for scale differences across methods. Fourier-MLP exhibits the lowest absolute and relative residuals, consistent with its strong smoothness bias. SVGP displays large absolute residuals despite low reconstruction error, indicating oversmoothing that fits sensor observations while deviating from the dynamics. FieldFormer achieves moderate residual values, balancing reconstruction accuracy with physical plausibility. These diagnostics reinforce the point that physics residuals should be interpreted as secondary consistency checks rather than primary indicators of knowledge recovery.

\section{Additional Results}

Table~\ref{tab:ablation_real_holdout} reports sensor-holdout results, used as a proxy for spatial generalization when dense full-field ground truth is unavailable. FieldFormer is stronger on most atmospheric holdout variables, while Gamma Field improves PM$_{10}$ holdout performance on the pollution dataset. The mixed results indicate that local coordinate-dependent scaling is useful in some heterogeneous regimes, but is not a universal substitute for observational support at unseen sensors.

\begin{table}[htbp]
\centering
\caption{Ablation study on real-world benchmarks: sensor-holdout evaluation. Results are reported as mean $\pm$ std over bootstrap samples. Lower is better.}
\label{tab:ablation_real_holdout}
\resizebox{0.95\textwidth}{!}{%
\begin{tabular}{llcccc}
\hline
Dataset & Variable 
& FieldFormer RMSE & FieldFormer MAE
& Gamma Field RMSE & Gamma Field MAE \\
\hline

Atmospheric & AT
& \textbf{4.63 $\pm$ 0.6376}
& \textbf{3.14 $\pm$ 0.4384}
& 4.683 $\pm$ 0.6321
& 3.201 $\pm$ 0.46 \\

Atmospheric & RH
& 15.75 $\pm$ 2.517
& \textbf{10.04 $\pm$ 1.568}
& \textbf{15.63 $\pm$ 1.952}
& 10.11 $\pm$ 1.318 \\

Atmospheric & $U_x$
& \textbf{2.26 $\pm$ 0.6081}
& \textbf{1.282 $\pm$ 0.2269}
& 2.362 $\pm$ 0.6137
& 1.33 $\pm$ 0.2493 \\

Atmospheric & $V_y$
& \textbf{1.502 $\pm$ 0.1426}
& \textbf{1.026 $\pm$ 0.09529}
& 1.671 $\pm$ 0.1498
& 1.122 $\pm$ 0.1189 \\

Pollution & PM$_{10}$
& 115 $\pm$ 13.11
& 69.63 $\pm$ 9.17
& \textbf{92.63 $\pm$ 9.512}
& \textbf{56.73 $\pm$ 6.764} \\

Pollution & PM$_{2.5}$
& \textbf{45.95 $\pm$ 4.77}
& \textbf{26.76 $\pm$ 3.152}
& 47.5 $\pm$ 4.847
& 27.14 $\pm$ 3.161 \\

\hline
\end{tabular}%
}
\end{table}
